\documentclass[12pt,a4paper]{article}

% Packages
\usepackage[utf8]{inputenc}
\usepackage[T1]{fontenc}
\usepackage{amsmath}
\usepackage{amssymb}
\usepackage{amsthm}
\usepackage{graphicx}
\usepackage{hyperref}
\usepackage{listings}
\usepackage{xcolor}
\usepackage{geometry}

% Geometry settings
\geometry{margin=1in}

% Theorem environments
\newtheorem{theorem}{Theorem}[section]
\newtheorem{lemma}[theorem]{Lemma}
\newtheorem{proposition}[theorem]{Proposition}
\newtheorem{corollary}[theorem]{Corollary}

\theoremstyle{definition}
\newtheorem{definition}{Definition}[section]
\newtheorem{example}{Example}[section]

\theoremstyle{remark}
\newtheorem*{remark}{Remark}
\newtheorem*{note}{Note}

% Custom commands
\newcommand{\R}{\mathbb{R}}
\newcommand{\N}{\mathbb{N}}
\newcommand{\Z}{\mathbb{Z}}
\newcommand{\Q}{\mathbb{Q}}
\newcommand{\C}{\mathbb{C}}

% Code listing settings
\lstset{
    language=Python,
    basicstyle=\ttfamily\small,
    keywordstyle=\color{blue},
    commentstyle=\color{green!60!black},
    stringstyle=\color{red},
    numbers=left,
    numberstyle=\tiny,
    breaklines=true,
    frame=single
}

% Title information
\title{One Monokai Python Theme: \\ LaTeX Sample Document}
\author{Theme Testing}
\date{\today}

\begin{document}

\maketitle

\begin{abstract}
This is a sample \LaTeX{} document to test the \textbf{One Monokai Python} theme's syntax highlighting for \LaTeX{} files. It includes various mathematical expressions, theorems, code listings, and common \LaTeX{} constructs.
\end{abstract}

\tableofcontents

\section{Introduction}

\LaTeX{} is a high-quality typesetting system that includes features designed for the production of technical and scientific documentation. This document demonstrates various \LaTeX{} features.

\subsection{Basic Text Formatting}

Here we have \textit{italic text}, \textbf{bold text}, \texttt{monospaced text}, and \underline{underlined text}. We can also use \emph{emphasized text} which adapts to context.

\section{Mathematical Expressions}

\subsection{Inline Mathematics}

The quadratic formula is given by $x = \frac{-b \pm \sqrt{b^2 - 4ac}}{2a}$. Euler's identity states that $e^{i\pi} + 1 = 0$.

\subsection{Display Mathematics}

The integral of a function $f(x)$ from $a$ to $b$ is denoted as:
\[
\int_a^b f(x) \, dx
\]

The sum of the first $n$ natural numbers:
\begin{equation}
\sum_{i=1}^{n} i = \frac{n(n+1)}{2}
\label{eq:sum}
\end{equation}

\subsection{Matrix Notation}

A $2 \times 2$ matrix can be written as:
\[
A = \begin{pmatrix}
a & b \\
c & d
\end{pmatrix}
\]

The identity matrix:
\begin{equation}
I_n = \begin{bmatrix}
1 & 0 & \cdots & 0 \\
0 & 1 & \cdots & 0 \\
\vdots & \vdots & \ddots & \vdots \\
0 & 0 & \cdots & 1
\end{bmatrix}
\end{equation}

\subsection{Aligned Equations}

\begin{align}
f(x) &= x^2 + 2x + 1 \\
     &= (x + 1)^2 \\
     &= (x + 1)(x + 1)
\end{align}

\subsection{Cases and Piecewise Functions}

\begin{equation}
f(x) = \begin{cases}
x^2 & \text{if } x \geq 0 \\
-x^2 & \text{if } x < 0
\end{cases}
\end{equation}

\section{Theorems and Proofs}

\begin{theorem}[Pythagorean Theorem]
\label{thm:pythagoras}
In a right triangle, the square of the hypotenuse is equal to the sum of squares of the other two sides:
\[
a^2 + b^2 = c^2
\]
\end{theorem}

\begin{proof}
The proof follows from the geometric interpretation... (proof omitted for brevity).
\end{proof}

\begin{lemma}
\label{lem:example}
For all $n \in \N$, we have $n^2 \geq n$.
\end{lemma}

\begin{definition}[Continuous Function]
A function $f: \R \to \R$ is continuous at $x_0$ if for every $\epsilon > 0$, there exists $\delta > 0$ such that:
\[
|x - x_0| < \delta \implies |f(x) - f(x_0)| < \epsilon
\]
\end{definition}

\begin{example}
Consider the function $f(x) = x^2$. This function is continuous everywhere on $\R$.
\end{example}

\section{Lists}

\subsection{Itemized List}

\begin{itemize}
    \item First item
    \item Second item
    \begin{itemize}
        \item Nested item A
        \item Nested item B
    \end{itemize}
    \item Third item
\end{itemize}

\subsection{Enumerated List}

\begin{enumerate}
    \item First step
    \item Second step
    \item Third step
    \begin{enumerate}
        \item Sub-step 3.1
        \item Sub-step 3.2
    \end{enumerate}
    \item Fourth step
\end{enumerate}

\section{Tables and Figures}

\subsection{Tables}

\begin{table}[h]
\centering
\begin{tabular}{|c|c|c|}
\hline
\textbf{Name} & \textbf{Value} & \textbf{Unit} \\
\hline
Speed of light & $3 \times 10^8$ & m/s \\
Planck constant & $6.626 \times 10^{-34}$ & J·s \\
Avogadro constant & $6.022 \times 10^{23}$ & mol$^{-1}$ \\
\hline
\end{tabular}
\caption{Physical constants}
\label{tab:constants}
\end{table}

\subsection{Figures}

\begin{figure}[h]
\centering
% \includegraphics[width=0.5\textwidth]{example-image}
\caption{Example figure (image not included)}
\label{fig:example}
\end{figure}

\section{Code Listings}

Here's an example of embedded Python code:

\begin{lstlisting}[caption={Python Fibonacci Implementation}]
def fibonacci(n):
    """Calculate Fibonacci sequence."""
    if n <= 0:
        return []
    elif n == 1:
        return [0]
    
    fib = [0, 1]
    for i in range(2, n):
        fib.append(fib[i-1] + fib[i-2])
    
    return fib

# Test the function
result = fibonacci(10)
print(f"Fibonacci sequence: {result}")
\end{lstlisting}

\section{Cross-References}

We can reference Equation~\ref{eq:sum} from Section~\ref{sec:math}. Theorem~\ref{thm:pythagoras} is a fundamental result in geometry. See Table~\ref{tab:constants} for physical constants.

\section{Special Symbols and Greek Letters}

Greek letters: $\alpha, \beta, \gamma, \delta, \epsilon, \zeta, \eta, \theta, \lambda, \mu, \pi, \rho, \sigma, \tau, \phi, \chi, \psi, \omega$

Capital Greek letters: $\Gamma, \Delta, \Theta, \Lambda, \Xi, \Pi, \Sigma, \Phi, \Psi, \Omega$

Mathematical operators: $\pm, \times, \div, \cdot, \leq, \geq, \neq, \approx, \equiv, \subset, \subseteq, \in, \notin, \cup, \cap, \emptyset$

Calculus: $\partial, \nabla, \infty, \lim, \int, \oint, \sum, \prod$

\section{Advanced Mathematics}

\subsection{Limits}

\[
\lim_{n \to \infty} \left(1 + \frac{1}{n}\right)^n = e
\]

\subsection{Derivatives}

\[
\frac{d}{dx} \left( \sin(x) \right) = \cos(x)
\]

\subsection{Integrals}

\[
\int_{-\infty}^{\infty} e^{-x^2} \, dx = \sqrt{\pi}
\]

\subsection{Partial Derivatives}

\[
\frac{\partial^2 u}{\partial t^2} = c^2 \nabla^2 u
\]

\section{Conclusion}

This document demonstrates various \LaTeX{} features for testing the One Monokai Python theme. The theme should properly highlight commands, environments, math mode, comments, and other \LaTeX{} constructs.

\begin{thebibliography}{99}

\bibitem{lamport1994}
Lamport, L. (1994).
\textit{\LaTeX: A Document Preparation System}.
Addison-Wesley Professional.

\bibitem{knuth1986}
Knuth, D. E. (1986).
\textit{The \TeX book}.
Addison-Wesley Professional.

\end{thebibliography}

\appendix

\section{Appendix: Additional Formulas}

\subsection{Taylor Series}

\[
f(x) = f(a) + f'(a)(x-a) + \frac{f''(a)}{2!}(x-a)^2 + \frac{f'''(a)}{3!}(x-a)^3 + \cdots
\]

\subsection{Fourier Transform}

\[
\hat{f}(\omega) = \int_{-\infty}^{\infty} f(t) e^{-i\omega t} \, dt
\]

\end{document}
